These labels are heuristic physical summaries, not confirmed planet taxonomies.

\begin{figure}[H]
\centering
\includegraphics[width=0.95\linewidth]{figures/01_main_graphs_by_population.pdf}
\caption{Grafos principales coloreados por poblacion heuristica.}
\end{figure}

La Figura 4 muestra una regularidad importante: casi todos los espacios contienen un componente central muy dominante y una periferia de nodos pequenos. Esa geometria sugiere que el catalogo tiene un ``bulk'' poblacional grande y varias subregiones mas especializadas en los extremos del espacio de variables. En el grafo orbital, la periferia es mas rica y mas extendida que en varios espacios fisicos, lo que respalda la idea de que la estructura orbital separa mejor regimenes de periodo y distancia estelar.

\begin{figure}[H]
\centering
\includegraphics[width=0.95\linewidth]{figures/02_main_graphs_by_imputation.pdf}
\caption{Grafos principales coloreados por imputacion.}
\end{figure}

La Figura 5 explica por que no basta con mirar forma topologica. En orbital y en los espacios fisicos, el componente central y buena parte de la periferia conservan colores compatibles con baja imputacion, lo que hace mas defendible su lectura fisica. En cambio, el espacio termico presenta anillos y regiones grandes dominadas por colores de alta imputacion; eso indica que parte de su complejidad visual esta construida sobre valores completados y no sobre observaciones directas. El espacio conjunto queda en una situacion intermedia: interesante, pero con mas dependencia de completacion que orbital.

La tabla de nodos destacados debe leerse como un primer mapa de regiones candidatas, no como una taxonomia final. Lo valioso aqui es que ya aparecen nodos compatibles con sub-Neptunos, supertierras y gigantes de periodo largo, junto con una indicacion explicita de cautela para nodos donde la estructura puede depender mas del pipeline que de una separacion fisica robusta.

\input{tables/highlighted_nodes_summary.tex}
